\chapter{Related Work}
\label{sec:rel}
\section{Background}
Business Process Management (BPM) refers to the systematic design, analysis, and optimization of organizational processes to improve efficiency and effectiveness, while Business Process Modeling (BPM) focuses on representing these processes in a structured and understandable form. BPM languages can be broadly categorized into conceptual, formal, and executable types. Conceptual languages, such as Business Process Model and Notation (BPMN), provide intuitive graphical notations that facilitate communication between business and technical stakeholders, although they lack fully formal semantics \cite{omg2011bpmn}. In contrast, formal languages like Petri Net offer mathematically precise representations, enabling rigorous analysis of properties such as correctness and deadlock-freedom \cite{vanderAalst1998}. To operationalize process models, executable structures such as Refined Process Structure Tree (RPST) are used to transform models into well-structured, executable components \cite{vanhatalo2009rpst}.

The generation of process models from textual descriptions, commonly referred to as text-to-model approaches, has evolved significantly over time. Early work relied on rule-based methods in Natural Language Processing (NLP), mapping linguistic elements such as verbs and nouns to process activities and actors \cite{friedrich2011process}. Subsequent statistical and machine learning approaches improved robustness but required annotated data and extensive feature engineering. More recently, advances in Large Language Models (LLMs) have enabled new approaches to process model generation from natural language. Recent studies indicate that LLM-based approaches have the potential to streamline modeling tasks and reduce manual effort; however, they still require careful evaluation of the generated models \cite{daclin2024gai4bm, kourani2024llm}.

\section{AI-Assisted Process Modeling Tools}
% While many tools exist that deal with the analysis of large document sets, to the best of our knowledge, no integrated tool with the dedicated purpose of creating process models from document sets exists. Therefore, we explore related tools that support AI-assisted process modeling. In total, eleven AI-assisted process modeling tools are identified through the literature review (see Table \ref{tab:lr_search}). In the following, we briefly describe these tools and focus on their underlying approaches and key features.


While many tools exist that deal with the analysis of large document sets, to the best of our knowledge, no integrated tool with the dedicated purpose of creating process models from document sets exists. Therefore, we explore related tools that support AI-assisted process modeling.

In total, eleven AI-assisted process modeling tools are identified through the literature review (see Table \ref{tab:lr_search}). In the following, we describe these tools with a focus on their underlying approaches and, in particular, their distinguishing functionalities.

\begin{table}[htb!]
    \centering
    \caption{Literature review: horizontal research}
    \label{tab:lr_search}
    \floatfoot{Source: https://scholar.google.com/}
    \begin{tabular}{lccl}
        \textit{Keywords for search} & \textit{hits} & \textit{selected} & \textit{Selection criteria}\\
        \hline
        allintitle: process modeling generative AI & 12 & 1 & Process modeling tool\\
        allintitle: process modeling LLM[s] & 16+7 & 5+0 & Process modeling tool\\ 
        allintitle: process modeling large language & 49 & 1 & Process modeling tool\\
        allintitle: conversational process modeling & 8 & 1 & Process modeling tool\\
        allintitle: BPMN generative AI & 2 & 0 & Process modeling tool\\ 
        allintitle: BPMN LLM[s]& 12+8 & 2+1 & Process modeling tool\\
        allintitle: BPMN large language& 7 & 1 &Process modeling tool\\
        \textit{Results horizontal search} & 121 & 11 & \\
        \textit{Results vertical/backward search} &  & & \\
        Added papers &  & 1& Process modeling tool \\
        \textit{Overall:} &  & 12 & \\
        \hline
    \end{tabular}
\end{table}


In the 12 hits of the first search query ("allintitle: process modeling generative AI"), the tool ProMoAI \cite{kourani2024promoai} is identified. It is the first tool shown in the academic literature to leverage LLMs for automated
process model generation, introduced in 2024, and is frequently referenced in subsequent work.
ProMoAI enables the transformation of natural language into process models and supports iterative
refinement through feedback loops, demonstrating the feasibility of LLM-based process modeling.

In the second query ("allintitle: process modeling LLM"), several tools are identified, including BPMN-Chatbot \cite{kopke2024introducing}, BPMN Assist \cite{licardo2026bpmn}, KICoPro \cite{lauer2026human}, and HyperMode \cite{pardo2025direct}. BPMN-Chatbot differs from ProMoAI by focusing on token efficiency, using lightweight prompting strategies to significantly reduce computational cost while maintaining correctness. BPMN Assist distinguishes itself by introducing a structured intermediate representation (JSON) together with function-based editing, enabling incremental updates instead of full model regeneration. KICoPro takes a different direction by explicitly incorporating human-in-the-loop interaction, allowing users to guide and adjust the model generation process continuously. HyperMode further extends interaction capabilities by combining LLM-based generation with a direct manipulation interface, enabling users to operate on specific diagram elements rather than relying solely on textual prompts.  

Additionally, the tool BPMNGen is firstly introduced in \cite{horner2025towards} and further developed in \cite{horner2026automatically}. In contrast to purely LLM-based generation approaches, BPMNGen integrates natural language processing (NLP) techniques to analyze and structure textual process descriptions before transforming them into BPMN models. This allows the tool to better handle linguistic structures and improve the extraction of process-relevant information.

In the last search attempt combining "process modeling" with "large language", the tool PRODIGY is identified in \cite{ziche2024llm4pm}. Unlike standalone generation tools, it integrates retrieval mechanisms to incorporate external knowledge, enabling more context-aware process model construction.

In the search query ("allintitle: conversational process modeling"), the tool AutoBPMN.AI \cite{klievtsova2025autobpmn} is found. Its distinguishing aspect lies in directly generating executable process representations, thereby reducing the gap between process modeling and execution compared to tools that focus only on diagram generation.

Finally, combining “BPMN” with "generative AI", "LLM[s]", or "large language" leads to three additional tools: BPMN-Chatbot++ \cite{safan2025bpmn}, LLM4BPMNGen \cite{costa2025llm4bpmngen}, and Nala2BPMN \cite{nour2024nala2bpmn}. BPMN-Chatbot++ extends the original chatbot approach by improving iterative refinement capabilities, making interaction more effective compared to earlier versions. LLM4BPMNGen differs from purely generative approaches by emphasizing more controlled and structured generation, improving output consistency. Nala2BPMN distinguishes itself through a modular pipeline architecture that separates extraction and model construction, enabling better handling of complex and unstructured input compared to end-to-end generation methods.

Overall, the reviewed tools are predominantly based on conversational, prompt-driven approaches for transforming natural language descriptions into process models. A common characteristic is the use of iterative refinement, where models are progressively improved through repeated user interaction. In terms of input, text-based descriptions represent the dominant modality, while only a limited number of tools extend this to additional forms such as voice or audio. Regarding output, the generated results are primarily expressed in BPMN. Table \ref{tab:lr_tools} summarizes these tools with a focus on their distinguishing characteristics.

\begin{table}[htb!]
\centering
\caption{Overview of AI-assisted process modeling tools}
\label{tab:lr_tools}
\begin{tabular}{p{2.5cm} l p{2cm} p{2cm} p{4cm}}
\hline
\textbf{Tool} & \textbf{Year} & \textbf{Input} & \textbf{Output} & \textbf{Key Contribution} \\
\hline

ProMoAI & 2024 
& Text 
& BPMN/PNML
& Automated generation and iterative refinement \\
\hline
BPMN-Chatbot & 2024 
& Text/voice 
& BPMN
& Optimized for token efficiency \\
\hline
Nala2BPMN & 2024 
& Text 
& BPMN
& Hybrid LLM + algorithm pipeline \\
\hline
PRODIGY & 2024 
& Text 
& BPMN (text)
& RAG-based enterprise modeling \\
\hline
LLM4BPMNGen & 2025 
& Text/voice/audio
& BPMN
& Audio upload + structured prompting \\
\hline
BPMN-Chatbot++ & 2025 
& Text/voice 
& BPMN
& Improved prompting and interaction \\
\hline
AutoBPMN.AI & 2025 
& Text 
& CPEE Tree
& Executable model generation \\
\hline
HyperMode & 2025 
& Text
& BPMN 
& Direct manipulation + LLM interaction \\
\hline
BPMNGen & 2026 
& Text 
& BPMN
& Conversational generation + human-centered evaluation \\
\hline
BPMN Assist & 2026 
& Text 
& BPMN
& JSON-based intermediate for efficient editing \\
\hline
KICoPro & 2026 
& text
& BPMN
& Able to answer questions and explain descisions \\

\hline
\end{tabular}

\end{table}

% \section{AutoBPMN.AI}
% Here, we provide a more detailed in depth review of the AutoBPMN.AI tool, as we will leverage its LLM service for the automated process model generation in our proposed system.


% AutoBPMN.AI provides an LLM-backed service that turns natural language process descriptions into an \,\emph{executable} process representation (CPEE-Tree) \cite{klievtsova2025autobpmn}. The core interaction is conversational: the user describes a procedure in free-form text, and the service produces (and maintains) a structured model corresponding to that description.

% Model creation is performed iteratively over multiple turns. After the initial model draft is generated, the user can issue follow-up instructions (e.g., clarifications, corrections, additions, or restructuring requests). The service uses the conversational context to interpret each instruction and updates the executable model accordingly, enabling step-by-step refinement rather than a single one-shot generation.

% In summary, the AutoBPMN.AI LLM service operates as a dialogue-driven modeling loop with two key artifacts exchanged between user and system: natural language messages as input and an updated CPEE-Tree model as output at each iteration.

% In this thesis, we will leverage the LLM REST service provided by AutoBPMN.AI for model generation and refinement within our proposed system.